\subsection{P028 --- Semantic Web Technologies for Indoor Environmental Quality: A Review and Ontology Design}
\label{paper:P028}
\textbf{Authors:} Alex Donkers, Dujuan Yang, Bauke de Vries, Nico Baken\par
\textbf{Major Category:} BIM and Semantic Information\par
\textbf{Secondary Category:} Building Environment and IEQ, Reviews, Frameworks and Standards\par
\textbf{Keywords:} Semantic web, Linked data, Ontology, Indoor environmental quality, Building Performance Ontology, BIM, IoT\par
\paragraph{主な主張}
IEQ modellingに必要なBIM/IoT由来の異種情報を扱うため, building topologyとstatic/dynamic propertiesを統合するBuilding Performance Ontologyを構築し, 複数levels of detailを備えた拡張可能なsemantic data modelを提示した. \cite{Donkers:2022BOPIEQ}
\paragraph{新規性}
既存ontologyのsystematic reviewから共通design patternsを抽出し, topology, static properties, observations, actuationsをlightweight upper ontologyとlower-level ontologiesへ統合した点に新規性がある.
\paragraph{位置付け}
IEQ向けsemantic integrationの基盤研究であり, BIM, IoT, BMSの異種情報をbuilding contextと結び付けるontology設計として位置付ける.
